Elles me demandèrent de leur apprendre à tricoter : elles s’installaient autour de
moi, toujours pieds nus et assises sur leur talons, jamais sur des chaises. Elles étaient
fort habiles et confectionnaient par exemple de petites corbeilles tissées avec de la
laine pour y placer galettes et gâteaux au miel...

Le responsable du hameau, Pedro, était un ancien adjudant d’Algérie, parlant arabe,
qui organisait les chantiers pour les hommes et veillait à la bonne entente dans le
camp. Lui succédèrent Susini puis Normand. Quant aux enfants, deux classes, en
préfabriqué, avaient été construites dans la cour de l’ancienne école de Meyrueis :
elles ont subsisté jusqu’aux années 2000 ! Un ramassage scolaire avait été mis en
place et c’est René Fages -- Fafa -- (il avait été caché au Moulin d’Ayres après avoir été blessé
lors des combats de la Parade, 20 ans auparavant), qui faisait office de chauffeur.
A notre égard ils furent tous d’une gentillesse extrême, avec de multiples attentions,
de petits cadeaux et des bols de couscous suffisants pour une famille (très)
nombreuse ! ».

A Ferrussac, Mme. Richard-Valdeyron et son mari, me donne quelques précisions
complémentaires :

\begin{quote}
\textit{« Je leur apportais, peut-être une vingtaine de litres de lait mais ce n’était pas tous
les jours. Mon mari laissait le bidon au bord de la route mais c’est moi qui devais
l’apporter aux femmes ; souvent elles le laissaient au soleil, attendant qu’il caille en
devenant aigre.}

\textit{Certains avaient plusieurs femmes Il y en avait un qui en avait même trois C’était le
plus âgé, le caïd. Sa première femme, la plus âgée, avait tout le pouvoir, les deux
autres, bien plus jeunes, obéissaient. C’est lui qui faisait la circoncision pour les
garçons. Parfois le médecin (le Dr. Rouzier) venait vérifier, à cause de l’hygiène.}

\textit{C’était de bons clients pour les commerçants car ils payaient rubis sur l’ongle.. Si
l’appoint manquait, nul besoin de le demander le lendemain. Parfois les hommes
descendaient à Meyrueis ; ils allaient au café et, de temps à autre, c’est le taxi de
Combemale qui les ramenait, lorsqu’ils avaient trop bu. Il conduisait également, à
l’hôpital ceux qui étaient gravement malades. »}
\end{quote}

\section*{Les Harkis au travail}

Les harkis étaient employés par les \og Eaux et Forêts \fg{}. Ils formaient deux équipes :
La 1\textsuperscript{ère}. travaillait à la pépinière de Roquedols sous la direction de M. Boulet, la 2\textsuperscript{ème}
était affectée aux travaux en forêt : elle élargit la route empierrée entre le Villaret et
Rousses. Pour ce faire, elle démolit et fit disparaître les restes de l’ancien moulin de
Rousses à 100m en aval du hameau et aussi l’ancien porche de la maison Martin, (nos
cousins Martin de Ferrussac et Raffègues. La maison et la propriété Martin avaient été vendues vers
1960 aux Eaux et Forêts), par lequel passait le chemin. Trop étroit il empêchait le passage
des camionnettes ou camions pour le transport du bois.
Les harkis, continuèrent jusqu’à la Pépinière, et tracèrent la route forestière
jusqu’au col de la Caumette permettant une sortie directe de la vallée du Béthuzon vers
l’Aigoual.
